% !TeX program = pdflatex
\documentclass[
	themeColor=blue,
	lang=DE,
	titleStyle=2,
	jobExpStyle=2,
	eduStyle=2
]{../templates/cv_res}

\name{Brian Stephenson}
\nationality{Vereinigte Staaten}
\profilepicture{../res/beige_wall_cropped.jpg}
\githubLink{github.com/brianESte}

\begin{document}
\maketitle

\Section{Berufserfahrung}
\employmentEntry{Hilfswissenschaftler}{Autonomous Digitization Technologies bei Fraunhofer IGD}
	[Darmstadt]{03.2024 -- }{%
	{Entwicklung eines Algorithmus für Pfadplannung auf einem Objekt einfacher arbiträrer Geometrie},
	{Lektorierung von wissenschaftlichen Arbeiten von Mitarbeitern},
	{Umarbeitung von Code eines eingebetteten Projekts, um Programgröße anzupassen und die Schnittstelle zwischen Rechner und Mikrokonroller zu vereinheitlichen}
}
\employmentEntry{Hilfswissenschaftler WalkerChair Project}{SIM Institute, TU Darmstadt}
{04.2023 -- 09.2023}{%
	{Entwicklung eines LED-Kettentreiberbibliothek für die STM32 Plattform mit mehreren Anzeigefunktionen},
	{Verbesserung vom Schrittmotorverhalten und Erweiterung der zugehörigen Dokumentation},
	{Dokumentierung des elektrischen Systems und der Kabelung}
}
\employmentEntry*{Engineering Team Werkstudent}{sigo GmbH}[Darmstadt]{12.2019 -- 03.2021}{%
  {Erzeugung und Verwaltung von 3D-Modellen ihrer Fahrräder und Stationen},
  {Entwicklung von neuen Komponenten, um sich ändernde Kunden- und Betreueranforderungen zu erfüllen},
  {Erstellung einer detailierten Fahrradkonstruktionseinleitung, um Qualität und Reproduzierbarkeit zu erhöhen}
}

\Section{Bildung}
\educationEntry{Technische Universität Darmstadt}[Fachbereich Elektrotechnik und Informationstechnik]{10.2019 -- 11.2025}{M.Sc. Mechatronik Vertiefung: Robotik}{Note: 1,780}
\educationEntry*{Technische Universität Darmstadt}[Fachbereich Maschinenbau]{07.2015 -- 03.2017}{B.Sc. Maschinenbau}{Note: 2,020}

\Section{Ehrenamtliche Arbeit}
\employmentEntry*{Reparaturer}{Martinsviertel RepairCafé}[Darmstadt]{10.2021 -- }{%
	{Fehlerdiagnose und Reparatur von defekten elektronischen und mechanischen Geräte}
}

% \Section{Relevante Projekte}
% Cooperative Game Device - self designed PCBs using KiCad, with each module controlled by an STM32 microcontroller that communicate via CAN bus.

\Section{Kenntnisse}
\skillLevelList{%
	{Autodesk Inventor,2},
	{DS Solidworks,3},
	{Siemens NX,3},
	{MATLAB/Simulink,2}
}{circle}[black]
\hfil
\skillLevelList{%
	{FreeCAD,2},
	{LTSpice,1},
	{KiCad,4},
	{Autodesk Eagle,1},
}{circle}[black]
\hfil
\skillLevelList{%
	{LabVIEW,2},
	{C (langauge),2},
	{C++,3},
	{Python,3},
}{circle}[black]

\Section{Sprachen}
\begin{center}
	\large%
	\begin{tabular}{p{18mm} p{30mm} p{24mm} p{18mm} p{30mm}}
		\hline
		\textbf{Deutsch} & C1 & & \textbf{Englisch} & Muttersprache
	\end{tabular}
\end{center}

\end{document}
